\problemchapter{Progetto assegnato}

\section*{Contesto}

Le batterie agli ioni di litio sono la tecnologia di accumulo dominante
per veicoli elettrici e sistemi stazionari, ma il loro stato di salute
e di carica non è osservabile direttamente: va inferito da misure
indirette. Fra queste, la \emph{Galvanostatic Electrochemical Impedance
Spectroscopy} (\acrshort{geis}) è una delle più informative, perché la
curva di Nyquist dell'impedenza complessa a diverse frequenze contiene
firme distintive dei fenomeni elettrochimici interni alla cella.

Il progetto si inserisce nella linea di ricerca avviata dal gruppo del
Politecnico di Milano (\textit{THERMINIC 2025}) sulla
caratterizzazione \acrshort{geis} di celle pouch
\acrshort{licoo2}/grafite, in collaborazione con il quale è stato messo
a disposizione un dataset sperimentale di circa 200 curve di Nyquist
coprenti 5 livelli di invecchiamento (\texttt{Aging}), 5 stati di carica
(\texttt{SOC}) e 8 temperature.

\section*{Obiettivo}

Sviluppare e confrontare modelli di Machine Learning su tre task
complementari di interpretazione della curva di Nyquist, ciascuno
corrispondente a una diversa domanda predittiva di interesse
elettrochimico. Per ciascun task, il deliverable atteso è un
\textbf{notebook Jupyter} che documenti l'analisi, addestri i modelli,
ne valuti le prestazioni con una strategia di \emph{cross-validation}
coerente con il task e ne discuta i risultati.

\section*{Compiti}

Il lavoro è articolato in tre task indipendenti.

\paragraph{Task 1 --- Ricostruzione \emph{Leave-One-Out} di una curva
di Nyquist.}
Dato un sottoinsieme di curve di Nyquist disponibili, predire la curva
mancante \(\bigl(Z_{\text{real}}(\omega), Z_{\text{imag}}(\omega)\bigr)\)
completa per la coppia \((\textit{Aging}, \textit{Temperature})\)
esclusa. La validazione segue lo schema \emph{Leave-One-Group-Out} sui
gruppi \((\textit{Aging}, \textit{Temperature})\). Modelli da
confrontare: Ridge, Random Forest, Gradient Boosting, K-Nearest
Neighbors, Bagging.

\paragraph{Task 2 --- Classificazione \emph{Young vs Old}.}
A partire dalle 8 curve di Nyquist di una coppia
\((\textit{Aging}, \textit{SOC})\) tenuta interamente fuori dal
training, classificare la coppia come ``giovane''
(\texttt{Aging}~\(\le\)~2) o ``vecchia'' (\texttt{Aging}~\(>\)~2). La
validazione utilizza un \emph{hold-out} dell'intera coppia. Modelli da
confrontare: Random Forest, Gradient Boosting, \acrshort{svm} a kernel
\acrshort{rbf}, oltre a baseline lineari e \acrshort{knn}.

\paragraph{Task 3 --- Interpolazione di un livello di \texttt{Aging}.}
Predire l'intero livello di invecchiamento mancante a partire dagli
altri quattro livelli, mantenendo costanti \texttt{SOC} e
\texttt{Temperature}. La validazione segue lo schema
\emph{Leave-One-Group-Out} sul gruppo \texttt{Aging}. Modelli da
confrontare: Ridge, Random Forest, Gradient Boosting, K-Nearest
Neighbors, Bagging.

\section*{Vincoli e note}

Il progetto è svolto in coppia (Davide Corso, Marco Soldani) sotto la
supervisione del Prof.~Loris Cannelli. La parte scientifica
(generazione del dataset, modello fisico di riferimento) è opera del
gruppo del Politecnico di Milano; il presente lavoro reimplementa e
estende l'analisi originale dal punto di vista ingegneristico, senza
sostituire il contributo scientifico di riferimento.